\section{Execution Model}
We define how a program is executed.

\subsection{Program, Execution, and VM State}

A \textbf{program} $\PPP$ is a fixed RV+ binary, compiled from an Ethereum State Transition function.
Its inputs $\III$ are derived from the Ethereum state (\emph{initial state}), and
the execution witness $\WWW$,
and its outputs $\OOO$ are the final Ethereum state (\emph{final state}).

The program $\PPP$ operates on the \textbf{VM state} $S$, defined as concatenation of the following components:
\begin{itemize}
\item $\pc$: 32-bit program counter
\item $\regs[0..k-1]$: list of $k$ 32-bit register values
\item $\mem[]$: byte-addressed writable memory
\end{itemize}
We separately denote the code of the program by $\code$. It is immutable and read-only, and is not part of the VM state.

A single instruction step of the program execution is defined as a transition from
one VM state to another according to the semantics of the instruction $\op$ retrieved from
$\code$ at the program counter $\pc_1$. This is denoted as follows:
\begin{equation}\label{eq:op}
S:=(\pc_1,\regs_1,\mem_1)\xrightarrow{\op:=\code[\pc_1]} S':=(\pc_2,\regs_2,\mem_2)
\end{equation}
For example, if $\op$ is an addition instruction, which adds the first two registers and stores in the third register, then the semantics of the step is defined as
\begin{align}
\pc_2&=\pc_1+1,\\
\regs_2&=\regs_1\text{ except that } \regs_2[2]:=\regs_1[0]+\regs_1[1],\\
\mem_2&=\mem_1.
\end{align}
We will use the notation $(\pc_1,\regs_1,\mem_1)\rightarrow  (\pc_2,\regs_2,\mem_2)$ to denote that $S$ transitions to $S'$ by executing instruction
 $\op:=\code[\pc_1]$.


An \textbf{execution trace} for program $\code$ applied to inputs $\III$ is a sequence of VM states

\[
\TTT:\quad S_0 \rightarrow S_1 \rightarrow \cdots \rightarrow S_T
\]
where each step corresponds to execution of one instruction.


A \textbf{zkVM proof} for $(\PPP,\EE,\EE')$,
where $\EE$ and $\EE'$ are the initial and final Ethereum states, is a SNARK that proves the following:
$$
\begin{cases}
S_0\rightarrow\cdots\rightarrow S_T;\\
S_0:=(\pc_0:=0,\regs,\mem)\text{ where $\regs,\mem$ are derived from }\EE;\\
S_T:=(\pc_T,\regs',\mem')\text{ where $\regs',\mem'$ are derived from }\EE';\\
\end{cases}
$$

\subsection{Supported Instructions}
The zkVM supports the RV+ instruction set, which is an extension of the RV32IM instruction set.
Concretely, all operations:
\begin{itemize}
    \item conform to the state transition \eqref{eq:op}, i.e.
    transform $(\pc_1,\regs_1,\mem_1)$ to $(\pc_2,\regs_2,\mem_2)$ according to the semantics of the instruction $\op$ retrieved from $\code$ at the program counter $\pc_1$.
    \item are all standard RISC-V instructions\footnote{\url{https://riscv.org/technical/specifications/}} (executed according to the standard RISC-V semantics)
     or additional
    instructions (precompiles) for efficient execution of certain operations --- concretely,
    Keccak and Poseidon.
\end{itemize}
For our purpose the exact semantics of the instructions is not important,
and we treat them as black boxes with their own predicates modeling the correctness of the execution.

Concretely, for given instruction $op$, the predicate $\varphi_{op}$  has the semantics
\begin{multline}\label{eq:phiop}
\varphi_{op}(\pc_1,\regs_1,\mem_1,\pc_2,\regs_2,\mem_2)=\mathrm{TRUE} \;\Longleftrightarrow\;\\
 (\pc_1,\regs_1,\mem_1)\xrightarrow{\op} (\pc_2,\regs_2,\mem_2) \wedge \code[\pc_1]=op
\end{multline}

For instance, if $op$ is an ADD instruction, then $\mem_1=\mem_2$, $\pc_1+1=\pc_2$ and the registers change according to the ADD instruction.


 Exceptions are memory operations,
which we assume have the following syntax and semantics:
\begin{align}
\rea:\;&(\regs_1:=(addr,*,\ldots),\mem)\rightarrow (\regs_2:=(addr,x,\ldots),\mem),\notag\\
&\quad \mem[addr]=x,\label{eq:mem-op-read}\\
\writ:\;&(\regs_1:=(addr,x,\ldots),\mem)\rightarrow (\regs_2:=(addr,x,\ldots),\mem'),\notag\\
&\quad \mem'[i]=\begin{cases}
x & \text{if }i=addr,\\
\mem[i] & \text{otherwise}.
\end{cases}\label{eq:mem-op-write}
\end{align}

We decompose each memory predicate into a pc-and-register component
(independent of memory) and an explicit memory-access equation:
\begin{multline}
  \varphi_{\rea}(S_1,S_2) :=
    \varphi'_{\rea}(\pc_1,\regs_1,\pc_2,\regs_2)
    \;\wedge\; \mem_1[\regs_1[0]] = \regs_2[1]
    \;\wedge\; \mem_2 = \mem_1
  \label{eq:phi-read-decomp}
\end{multline}
\begin{multline}
  \varphi_{\writ}(S_1,S_2) :=
    \varphi'_{\writ}(\pc_1,\regs_1,\pc_2,\regs_2)
    \;\wedge\; \mem_2 = \mem_1[\regs_1[0]\mapsto\regs_1[1]]
  \label{eq:phi-write-decomp}
\end{multline}
where $\regs_1[0]$ denotes the address register, $\regs_1[1]$ (resp.\ $\regs_2[1]$)
the value register before (resp.\ after) the step, and
$\varphi'_{\rea}$, $\varphi'_{\writ}$ depend only on $\pc$ and $\regs$.
Both $\varphi'_{\rea}$ and $\varphi'_{\writ}$ include the well-formedness
conjunct $\regs_1[0]\in[n]$, where $n$ is the number of addressable memory
cells: without it the expressions $\mem_1[\regs_1[0]]$ and
$\mem_1[\regs_1[0]\mapsto\regs_1[1]]$ above would be undefined.  The conjunct is
memory-free, so it is inherited verbatim by the committed-memory versions of the
two predicates, and it is what lets the security proof treat $\regs_1[0]$ as a
legal index of the committed vector.

Every pc-and-register part $\varphi'_{op}$ --- of memory and non-memory
operations alike --- also carries the fetch conjunct $\code[\pc_1]=op$
of~\eqref{eq:phiop}; the conjunct depends only on $\pc_1$ (the code $\code$ is
fixed), so it is memory-free and belongs in $\varphi'_{op}$.  This placement is
load-bearing: the security analysis (Remark~\ref{rem:mem-ext-hyp4}) relies on
it to conclude that at most one operation's disjunct of the step predicate can
hold, and that this operation is $\code[\pc_1]$.
\begin{instruction}
Teams whose architecture selects the operation differently (e.g.\ via selector
columns) must ensure an equivalent guarantee: the constraint system must force
the satisfied operation branch to be the one fetched at $\pc_1$.  Without it,
the memory-reconstruction rule of
Proposition~\ref{prop:memory-extractability} can diverge from the satisfied
branch and the commitment invariant fails without any binding violation.
\end{instruction}

The conjunct $\mem_2=\mem_1$ in~\eqref{eq:phi-read-decomp} is kept explicit for
the same reason as for the non-memory operations: a read must not be able to
silently alter memory.  Without it $\varphi_{\rea}$ would leave $\mem_2$ completely
unconstrained, and a trace satisfying $\varphi_{\step}$ at every step could
still replace the whole memory at a read.  Note that $\varphi_{\writ}$ needs no
such conjunct, as~\eqref{eq:phi-write-decomp} already determines $\mem_2$
entirely.


